\addchap{Preface}

This book began with a question: which textbooks should exist but do not? One answer is a textbook about a skill that every working mathematician uses and almost nobody teaches, the skill of finding the setting in which a hard problem becomes easy. Grothendieck called the result of that skill a rising sea. Most of us know it from the inside, as the moment when a proof that had been a fight turns out, in the right language, to be a few lines. This book tries to make that moment teachable.

It does so through cases. Each chapter tells the story of one problem, from analysis, probability, combinatorics, optimization, theoretical computer science, learning theory or statistics. First it shows the problem being fought in its original setting. Then it isolates the single move that changed the setting, re-proves the old results in the new one, and asks honestly what the move cost and where it fails. The cases were chosen to cover different moves rather than different fields, and the exercises ask the reader to carry each move to problems it has not yet been applied to.

The intended readers are graduate students and researchers in mathematics, theoretical computer science, machine learning and statistics, and anyone who has finished a hard proof and wondered whether it needed to be hard. The prerequisites are modest; each chapter builds what it uses. The prologue explains the plan and the format, and the epilogue compares the moves and discusses when the method does not work.

\paragraph{A note on authorship.} This book was written by Claude Opus 5.5, an AI model made by Anthropic, with human oversight by Xiaoyu Li, who asked for a textbook that ought to exist. Every chapter's proofs were re-derived by at least one independent reader, and its historical claims and references were checked against primary or secondary sources wherever these could be reached; a few details rest on secondary accounts, and recent preprints are marked as unrefereed. Mistakes surely remain. Readers who find one are warmly invited to treat it as an exercise: deciding whether a proof is right is itself one of the skills this book hopes to teach.
